\chapter{Trace-free curvature and CMC pinching}\label{sec:cmc-pinching}
The trace-free second fundamental form measures the failure to be umbilic. We derive its equation in a space form and use a sharp algebraic bound to prove a curvature gap on closed hypersurfaces. The spherical comparison source is Alencar--do Carmo \cite{AdC94}, Theorem 1.2 and Section 2. Their shape operator has the opposite sign and their mean curvature is the average of the principal curvatures. Thus their scalar mean curvature is $-H/n$ in our convention; their tensor $\phi$ agrees with the trace-free tensor used below. Their classification of all equality cases is not used in our proofs.

\begin{definition}[Trace-free second fundamental form]\label{def:tracefree-second-form}
For a smooth two-sided hypersurface of dimension $n\geq2$, set
\[
 b=h-\frac Hn g,\qquad B=A-\frac Hn I,\qquad s=|B|.
\]
Then $\operatorname{tr}B=0$ and $|A|^2=s^2+H^2/n$. The hypersurface is totally umbilic exactly when $B=0$. The quantity $s$ does not depend on the choice of normal. Compare the tensor $\phi$ in \cite{AdC94}, Section 1.
\end{definition}

\begin{proposition}[The trace-free Simons equation for CMC hypersurfaces]\label{prop:cmc-tracefree-simons}
Let $\Sigma^n$ be a smooth two-sided immersed hypersurface with constant $H$ in a space form of sectional curvature $c$. With the preceding notation and $\alpha=nc+H^2/n$,
\[
 \Delta b=(\alpha-s^2)b+H\left(b^2-\frac{s^2}{n}g\right),
\]
\[
 \frac12\Delta s^2=|\nabla b|^2+\alpha s^2-s^4+H\operatorname{tr}(B^3).
\]
Here $(b^2)_{ij}=b_i{}^kb_{kj}$. For $c=1$ the second equation is the normalization of \cite{AdC94}, equation (2.4).
\end{proposition}
\begin{proof}
Since $H$ is constant and $\nabla g=0$, one has $\Delta b=\Delta h$. In Theorem \ref{thm:spaceform-simons}, substitute $h=b+(H/n)g$ and $|A|^2=s^2+H^2/n$. The coefficient of $b$ becomes
$2H^2/n-s^2-H^2/n+nc=\alpha-s^2$. The scalar multiple of $g$ becomes $-Hs^2/n$, while the quadratic term is $Hb^2$. This proves the tensor equation. Contract it with $b$, use $\langle b,g\rangle=0$, and apply $\frac12\Delta|b|^2=|\nabla b|^2+\langle b,\Delta b\rangle$ to obtain the norm equation.
\end{proof}

\begin{lemma}[The sharp cubic bound for a trace-free symmetric operator]\label{lem:tracefree-cubic-bound}
For a symmetric operator $B$ on an $n$-dimensional Euclidean space, $n\geq2$, with $\operatorname{tr}B=0$, put $\beta_n=(n-2)/\sqrt{n(n-1)}$. Then
\[
 |\operatorname{tr}(B^3)|\leq\beta_n|B|^3.
\]
For $n>2$ and $B\ne0$, equality holds exactly when one eigenvalue has multiplicity one and the other has multiplicity $n-1$. For $n=2$ both sides vanish for every such $B$. This is the algebraic inequality attributed to Okumura in \cite{AdC94}, Section 2, Lemma 1; we include the optimization argument.
\end{lemma}
\begin{proof}
The zero operator is immediate. Divide the eigenvalues by $|B|$; they now satisfy $\sum_i\lambda_i=0$ and $\sum_i\lambda_i^2=1$. For $n=2$ they are opposite, so their cubes cancel. Suppose $n>2$. The normalized constraint set is compact and its two constraint gradients are independent: otherwise all $\lambda_i$ would be equal, contrary to the constraints. At an extremum of $\sum_i\lambda_i^3$, Lagrange multipliers give $3\lambda_i^2=a+2b\lambda_i$. There are at most two distinct eigenvalues, and the constraints rule out a single value.

Let the positive value have multiplicity $p$ and the negative value multiplicity $q=n-p$. The constraints give
\[
 \lambda_+=\sqrt{\frac{q}{np}},\qquad
 \lambda_-=-\sqrt{\frac{p}{nq}},\qquad
 \sum_i\lambda_i^3=\frac{q-p}{\sqrt{npq}}.
\]
The square of the last expression is $n/(pq)-4/n$. Since $pq=p(n-p)\geq n-1$, its maximum is $(n-2)^2/(n(n-1))$, attained precisely when $p=1$ or $q=1$. This bounds both extrema and proves the inequality after restoring the factor $|B|^3$. It also proves the equality assertion, including its converse.
\end{proof}

\begin{theorem}[A pinching gap for closed CMC hypersurfaces]\label{thm:cmc-pinching-gap}
Let $\Sigma^n$ be a smooth connected compact two-sided immersed hypersurface without boundary in a space form of curvature $c$, with constant scalar mean curvature $H$. Assume $n\geq2$ and $\alpha=nc+H^2/n>0$. Define
\[
 r=\frac{\sqrt{\beta_n^2H^2+4\alpha}-\beta_n|H|}{2}>0,
 \qquad \beta_n=\frac{n-2}{\sqrt{n(n-1)}}.
\]
If $|B|\leq r$ everywhere, then $B$ is parallel and either $B=0$ everywhere or $|B|=r$ everywhere. In particular, $\sup_\Sigma|B|<r$ forces the hypersurface to be totally umbilic. For $c=1$ this is the gap part of \cite{AdC94}, Theorem 1.2(i), in our normalization; the proof below also covers the stated values of $c$ with $\alpha>0$.
\end{theorem}
\begin{proof}
Proposition \ref{prop:cmc-tracefree-simons} and Lemma \ref{lem:tracefree-cubic-bound} imply
\[
 \frac12\Delta s^2\geq|\nabla b|^2+s^2(\alpha-s^2-\beta_n|H|s).
\]
The polynomial $P(s)=\alpha-s^2-\beta_n|H|s$ is positive on $[0,r)$ and zero at $r$. Under the pinching hypothesis the right side is nonnegative. Integrating over the compact manifold without boundary gives zero on the left. The two continuous nonnegative terms on the right must therefore vanish everywhere. Thus $\nabla b=0$, so connectedness makes $s$ constant, and $s^2P(s)=0$ forces $s=0$ or $s=r$. The strict bound excludes the latter. Compactness and absence of boundary justify this integration; the same conclusion has not been asserted for arbitrary noncompact pieces of CMC hypersurfaces.
\end{proof}

\begin{corollary}[The first curvature gap for a spherical minimal hypersurface]\label{cor:minimal-sphere-curvature-gap}
Let $\Sigma^n$ be a smooth closed connected two-sided immersed minimal hypersurface in the unit sphere, $n\geq2$. If $|A|^2\leq n$, then $A$ is parallel and either $A=0$ everywhere or $|A|^2=n$ everywhere. In particular, $|A|^2<n$ everywhere implies that the hypersurface is totally geodesic. This is the gap assertion recalled in \cite{AdC94}, Theorem 1.1(i), attributed there to Simons.
\end{corollary}
\begin{proof}
Here $H=0$, $c=1$, $B=A$, $\alpha=n$, and $r=\sqrt n$. Apply Theorem \ref{thm:cmc-pinching-gap}. In the strict case compactness also gives a strict supremum bound; alternatively the constant-norm conclusion excludes the value $n$. Vanishing of the second fundamental form is exactly the totally geodesic condition. This conclusion does not require the classification of the nonzero equality case.
\end{proof}

\begin{example}[Products of spheres at and beyond the pinching threshold]\label{ex:cmc-product-pinching}
For $p,q\geq1$, $n=p+q$, and $a,b>0$ with $a^2+b^2=1$, the product
\[
 \Sigma=\mathbb S^p(a)\times\mathbb S^q(b)\subset\mathbb S^{n+1}(1)
\]
is a closed embedded CMC hypersurface with parallel second fundamental form. For a suitable normal its principal curvatures are $b/a$ with multiplicity $p$ and $-a/b$ with multiplicity $q$, and
\[
 H=\frac{pb}{a}-\frac{qa}{b},\qquad
 |B|^2=\frac{pq}{n a^2b^2}.
\]
It is minimal precisely when $a^2=p/n$, in which case $|A|^2=n$. For $n=2$ every member realizes equality in Theorem \ref{thm:cmc-pinching-gap}. For $n>2$, $H\ne0$, equality requires one of $p,q$ to be one and additionally $H(q-p)<0$. This illustrates why not every nonminimal product is at the first pinching threshold. Compare \cite{AdC94}, Section 1 and Theorem 1.2.
\end{example}
\begin{proof}
At $(x,y)$ the vector $\nu=((b/a)x,-(a/b)y)$ is unit, orthogonal to $(x,y)$, and normal to the product within the sphere. Differentiating in the two factor directions gives the displayed principal curvatures. The splitting of the product tangent bundle is parallel, and these principal curvatures are constant; hence $A$ is parallel. Their trace is $H$. For two eigenvalues $k_+,k_-$ of multiplicities $p,q$, the squared trace-free norm is $(pq/n)(k_+-k_-)^2$. Here $k_+-k_-=1/(ab)$, giving the stated formula. Solving $H=0$ yields $a^2=p/n$, and substitution gives $|A|^2=n$.

For $n=2$, the identity $|B|^2=2+H^2/2$ follows either from the displayed formulas or from $k_+k_-=-1$. This is $r^2$ since $\beta_2=0$. For $n>2$, the nonzero trace-free eigenvalues are $q/(nab)$ and $-p/(nab)$, so $\operatorname{tr}(B^3)$ has sign $q-p$. The parallel norm identity gives
$0=\alpha s^2-s^4+H\operatorname{tr}(B^3)$.
At $s=r$, the inequality used in the gap theorem is an equality, requiring $H\operatorname{tr}(B^3)=-|H|\beta_n s^3$. For $H\ne0$, the sharp cubic equality requires $p=1$ or $q=1$, and its sign requires $H(q-p)<0$. Conversely these two conditions make the norm identity read $s^2(\alpha-s^2-\beta_n|H|s)=0$; since $s>0$, its unique positive root is $r$. No classification of other hypersurfaces is used.
\end{proof}
